\chapter{The Convex Quadratic Program and Its Certificate}
\label{ch:certificate}

This chapter proves the four facts everything later rests on: that~\eqref{eq:qp}
has exactly one solution; that the Karush--Kuhn--Tucker conditions characterise
it; that they are \emph{sufficient} and not merely necessary; and that the
working-set subproblem has a closed-form solution. The last section takes the
dual, which explains where the method of Chapter~\ref{ch:walking} gets its name
and why its objective increases.

The mathematics is short. The consequence---that a candidate answer can be
checked rather than trusted---is the organising idea of the book, and
Section~\ref{sec:certificate-principle} states it as a design principle rather
than a theorem.

\section{Existence and uniqueness}

Let
\[
  \Omega = \{x \in \R^n : c_i\T x = b_i \ (i \le m_{\mathrm{eq}}),\ \
                          c_i\T x \ge b_i \ (i > m_{\mathrm{eq}})\}
\]
denote the feasible set, a closed convex polyhedron, and recall
$f(x) = \tfrac12 x\T G x - a\T x$ with $G \succ 0$.

\begin{proposition}[Existence and uniqueness]
\label{prop:existunique}
If $\Omega \ne \emptyset$ then~\eqref{eq:qp} has exactly one solution.
\end{proposition}

\begin{proof}
Since $G \succ 0$, writing $\lambda_{\min} > 0$ for its smallest eigenvalue gives
$f(x) \ge \tfrac12\lambda_{\min}\norm{x}^2 - \norm{a}\norm{x} \to \infty$ as
$\norm{x}\to\infty$, so $f$ is coercive. Fix any $x_0 \in \Omega$; the sublevel set
$\{x \in \Omega : f(x) \le f(x_0)\}$ is then closed and bounded, hence compact and
non-empty, and a continuous function attains its minimum on it. That minimum is a
minimum over all of $\Omega$. Uniqueness: if $x \ne y$ both minimise, then by
strict convexity $f(\tfrac12(x+y)) < \tfrac12 f(x) + \tfrac12 f(y) = \min_\Omega f$,
and $\tfrac12(x+y) \in \Omega$ by convexity of $\Omega$, a contradiction.
\end{proof}

The proposition is unremarkable, but note which hypothesis did which work.
Coercivity came from $G \succ 0$ and gave existence; strict convexity came from
the same place and gave uniqueness. Drop to $G \succeq 0$ and both arguments fail,
though both conclusions may survive---a semidefinite QP on a compact polyhedron
still has a solution, and may have a whole face of them.

\section{Optimality conditions}

The Lagrangian of~\eqref{eq:qp} is
$L(x,\lambda) = \tfrac12 x\T G x - a\T x - \lambda\T(C\T x - b)$, and the
Karush--Kuhn--Tucker (KKT) conditions are
\index{Karush--Kuhn--Tucker conditions}
\begin{align}
  G x - a &= C \lambda, \tag{stationarity}\label{eq:stat}\\
  c_i\T x &= b_i \ (i \le m_{\mathrm{eq}}), \qquad
  c_i\T x \ge b_i \ (i > m_{\mathrm{eq}}), \tag{primal feasibility}\label{eq:pfeas}\\
  \lambda_i &\ge 0 \quad (i > m_{\mathrm{eq}}), \tag{dual feasibility}\label{eq:dfeas}\\
  \lambda_i\,(c_i\T x - b_i) &= 0 \quad (i > m_{\mathrm{eq}}).
      \tag{complementarity}\label{eq:comp}
\end{align}

Multipliers of equality constraints are unrestricted in sign. That single
asymmetry propagates into every formula in Part~II, and forgetting it is the most
common implementation error in this subject.

The four conditions have a reading worth internalising before any algebra.
Stationarity says the objective's gradient at $x$ is a combination of the
constraint normals: there is no direction of descent that the constraints permit.
Dual feasibility says the combination uses the inequality normals with the right
sign, pushing \emph{into} the feasible set rather than out of it. Complementarity
says only constraints that actually touch $x$ get to contribute. Primal
feasibility says $x$ is in $\Omega$ at all.

That the conditions are \emph{necessary} at a solution follows from any standard
constraint qualification, and for linear constraints no qualification is needed:
linearity is itself enough (see \cite[Ch.~12]{nocedal2006}). What matters here is
the converse.

\section{Sufficiency: the certificate}
\label{sec:certificate-principle}

\begin{proposition}[Sufficiency]
\label{prop:sufficient}
\index{sufficiency of KKT}
If $(x,\lambda)$ satisfies \eqref{eq:stat}--\eqref{eq:comp} then $x$ is the unique
solution of~\eqref{eq:qp}.
\end{proposition}

\begin{proof}
Let $\tilde x \in \Omega$. By strict convexity,
$f(\tilde x) \ge f(x) + (G x - a)\T(\tilde x - x)$, with equality only if
$\tilde x = x$. Substituting~\eqref{eq:stat},
\[
  (G x - a)\T(\tilde x - x)
  = \lambda\T\bigl(C\T \tilde x - C\T x\bigr)
  = \sum_i \lambda_i\bigl[(c_i\T\tilde x - b_i) - (c_i\T x - b_i)\bigr]
  = \sum_i \lambda_i (c_i\T \tilde x - b_i),
\]
the last step by~\eqref{eq:comp} for the inequalities and by $c_i\T x = b_i$ for
the equalities. Each surviving term is non-negative: for
$i \le m_{\mathrm{eq}}$ the bracket vanishes, and for $i > m_{\mathrm{eq}}$ both
factors are non-negative by~\eqref{eq:pfeas} and~\eqref{eq:dfeas}. Hence
$f(\tilde x) \ge f(x)$ with equality only at $\tilde x = x$.
\end{proof}

The proof is four lines and it licenses most of the interesting algorithms in this
book. Stated as a design principle:

\begin{quote}
\textbf{The certificate principle.} A method that produces candidate pairs
$(x,\lambda)$ may be as unprincipled as one likes---a guess, a cached answer from
a neighbouring problem, the output of a heuristic---provided every candidate is
tested against \eqref{eq:stat}--\eqref{eq:comp}. Sufficiency is what makes an
unguaranteed method \emph{safe} rather than \emph{usually right}.
\end{quote}

In floating point the four conditions become four tolerances. A candidate is
accepted when
\begin{equation}
  \norm{G x - a - C\lambda}_\infty \le \varepsilon,
  \quad |c_i\T x - b_i| \le \varepsilon \ (i \le m_{\mathrm{eq}}),
  \quad c_i\T x - b_i \ge -\varepsilon, \ \ \lambda_i \ge -\varepsilon,
  \ \ |\lambda_i (c_i\T x - b_i)| \le \varepsilon
  \label{eq:certificate}
\end{equation}
on the inequalities, with $\varepsilon$ scaled by the data. Two points about this
test recur throughout the book and are worth flagging now.

First, \emph{stationarity is checked against $G$ directly}, not inherited from
whatever construction produced $x$. The construction is exactly what an
ill-conditioned working set corrupts, so trusting it defeats the purpose of
checking.

Second, the test is not delicate, because the quantity it thresholds is
\emph{bimodal}. A candidate built from a well-conditioned working set satisfies
\eqref{eq:certificate} to a few multiples of the machine epsilon; one built on a
set that is not has no reason to come anywhere close. The two populations lie
orders of magnitude apart and the threshold sits in the empty space between them.
Chapter~\ref{ch:numerics} returns to this pattern, which is the single most useful
thing to know about tolerances in this subject.

\section{The working-set subproblem}

Fix a working set $\Active$ of size $k$ with normals $A = C_{\cdot\Active}$ and
suppose $A$ has full column rank $k$. Recall the subproblem~\eqref{eq:sub}: minimise
$f$ subject to $A\T x = b_\Active$, with every constraint in $\Active$ held as an
equality regardless of its type in the original program and every constraint
outside $\Active$ ignored.

\begin{proposition}[Solution of the subproblem]
\label{prop:sub}
\index{working set!subproblem}
Problem~\eqref{eq:sub} has the unique solution and multipliers
\begin{equation}
  u = \bigl(A\T G^{-1} A\bigr)^{-1}\bigl(b_\Active - A\T x_u\bigr),
  \qquad
  x = x_u + G^{-1} A\, u,
  \qquad x_u := G^{-1} a .
  \label{eq:subsol}
\end{equation}
\end{proposition}

\begin{proof}
Stationarity for~\eqref{eq:sub} reads $Gx - a = Au$, so
$x = G^{-1}a + G^{-1}Au = x_u + G^{-1}Au$. Substituting into $A\T x = b_\Active$
gives $A\T x_u + (A\T G^{-1}A)u = b_\Active$. The matrix $A\T G^{-1}A$ is positive
definite because $G^{-1} \succ 0$ and $A$ has full column rank, hence invertible,
which yields~\eqref{eq:subsol}; the point so constructed is the unique minimiser by
strict convexity of $f$ on the affine feasible set of~\eqref{eq:sub}.
\end{proof}

\begin{definition}[Dual Hessian]
\label{def:dualhessian}
\index{dual Hessian}
The matrix $A\T G^{-1} A$ is the \emph{dual Hessian} of the working set $\Active$.
\end{definition}

The dual Hessian appears under three names in what follows: as the matrix whose
Cholesky factor the solver of Chapter~\ref{ch:walking} carries; as the coefficient
matrix of the linear system that the guessing method of Chapter~\ref{ch:guessing}
solves; and as the Gram matrix of a least-squares problem in
Lemma~\ref{lem:directions}. Its invertibility is equivalent to $A$ having full
column rank, and \emph{every rank condition in this book comes from that single
fact}.

The pair $(x,u)$ of~\eqref{eq:subsol}, extended by zeros off the working set,
satisfies stationarity~\eqref{eq:stat} and complementarity~\eqref{eq:comp} by
construction: the working-set constraints have zero slack and the rest have zero
multiplier. What it need \emph{not} satisfy is the sign condition~\eqref{eq:dfeas}
on the inequalities in $\Active$, or primal feasibility~\eqref{eq:pfeas} off
$\Active$. Two of the four conditions come free; the other two are what an
algorithm must work for. Which two an algorithm keeps and which it chases is
exactly what distinguishes a primal from a dual method.

\begin{remark}[Cold start]
\label{rem:cold}
For $\Active = \emptyset$ the subproblem is unconstrained and~\eqref{eq:subsol}
degenerates to $x = x_u = G^{-1}a$ with an empty multiplier vector. The sign
condition holds vacuously, so this state satisfies the requirements of a dual
method with no work beyond one Cholesky factorisation. The objective there is
$f(x_u) = -\tfrac12 a\T x_u$. \emph{A dual method's phase-one problem is nothing
more than this.}
\end{remark}

\begin{example}[Two variables, one constraint]
Take $n = 2$, $G = I$, $a = 0$, and the single inequality
$c_1 = (1,1)\T$, $b_1 = 1$. Then $x_u = 0$, which is infeasible. With
$\Active = \{1\}$, $A = c_1$, the dual Hessian is $c_1\T c_1 = 2$, so
$u = (1 - 0)/2 = 1/2$ and $x = 0 + I c_1 (1/2) = (1/2,1/2)\T$. The multiplier is
positive and the point is feasible, so by Proposition~\ref{prop:sufficient} this
is the solution: the closest point of the half-space to the origin, as it should
be. Note that the certificate was checked, not assumed.
\end{example}

\section{The dual problem}

The following explains the name of the method in Chapter~\ref{ch:walking} and,
more usefully, why its objective value can be carried as a running total and why
that total increases.

\begin{proposition}[Dual function]
\label{prop:dual}
\index{duality!Lagrangian dual of the QP}
The Lagrangian dual of~\eqref{eq:qp} is the concave program $\max\,\psi(\lambda)$
over $\lambda_i \ge 0$ $(i > m_{\mathrm{eq}})$, where
\begin{equation}
  \psi(\lambda) = -\tfrac12 (a + C\lambda)\T G^{-1} (a + C\lambda) + b\T\lambda .
  \label{eq:dualfn}
\end{equation}
If $(x,\lambda)$ satisfies stationarity~\eqref{eq:stat} and
complementarity~\eqref{eq:comp}, then $\psi(\lambda) = f(x)$.
\end{proposition}

\begin{proof}
Minimising $L(\cdot,\lambda)$ over $x$ gives $x(\lambda) = G^{-1}(a + C\lambda)$
and, writing $s = a + C\lambda$,
$\psi(\lambda) = \tfrac12 s\T G^{-1}s - s\T G^{-1}s + b\T\lambda$, which
is~\eqref{eq:dualfn}. For the second claim, \eqref{eq:stat} gives $s = Gx$, and
\eqref{eq:comp} together with $c_i\T x = b_i$ on the equalities gives
$b\T\lambda = (Gx - a)\T x$; adding the two yields $f(x)$.
\end{proof}

So along any sequence of states of the form~\eqref{eq:subsol} the tracked
objective value is simultaneously the primal objective at the subproblem
minimiser and the dual objective at the current multipliers. That gives the
sandwich
\[
  \psi(\lambda) \;\le\; f(x^\star) \;\le\; f(\tilde x)
  \qquad\text{for any dual feasible $\lambda$ and any primal feasible $\tilde x$,}
\]
so a method whose iterates are of the form~\eqref{eq:subsol} approaches the
optimal value \emph{from below}. It is the mirror image of a primal method, which
maintains feasibility and descends to $f(x^\star)$ from above. Everything about
the trade between the two follows from this picture: a primal method can be
stopped early and return a feasible suboptimal point; a dual method cannot, but
it needs no phase one.

\begin{notes}
Proposition~\ref{prop:sufficient} and the closed form~\eqref{eq:subsol} are
standard; the presentation here follows \cite{schmelzer2026quadprog}, which states
them in exactly this form because every certificate in its account of the dual method
rests on them. Nocedal and Wright \cite[Ch.~16]{nocedal2006} give
the general treatment of quadratic programming, and \cite[Ch.~12]{nocedal2006} the
constraint qualifications that the linearity of our constraints lets us skip.
Boyd and Vandenberghe \cite{boyd2004} develop Lagrangian duality in the generality
of which Proposition~\ref{prop:dual} is a special case.

The bimodality observation about tolerance~\eqref{eq:certificate} is empirical
rather than a theorem, and the measurements behind it are reported in the
experiments of \cite{schmelzer2026quadprog}; Chapter~\ref{ch:numerics} returns to it.
\end{notes}

\begin{exercises}

\begin{exercise}
Prove that $A\T G^{-1} A \succ 0$ if and only if $A$ has full column rank, given
$G \succ 0$. Where in Proposition~\ref{prop:sub} is each direction used?
\end{exercise}

\begin{exercise}
Verify Proposition~\ref{prop:dual} directly for $n = 1$, $G = 1$, $a = 0$,
$C = 1$, $b = 1$, $m_{\mathrm{eq}} = 0$: write $\psi$ explicitly, maximise it over
$\lambda \ge 0$, and check that the maximiser reproduces the primal solution
through $x = G^{-1}(a + C\lambda)$.
\end{exercise}

\begin{exercise}
Show that the certificate test~\eqref{eq:certificate} costs $O(n^2 + nm)$ for
dense $G$ and $C$, and identify which term dominates for the long-only portfolio
problem of Chapter~\ref{ch:oneproblem}, where $m \approx n$.
\end{exercise}

\begin{exercise}
Give an example with $G \succeq 0$ singular in which the KKT conditions hold at
two distinct points. Which step of the proof of Proposition~\ref{prop:sufficient}
fails?
\end{exercise}

\begin{exercise}
Suppose an algorithm returns $(x,\lambda)$ with $\lambda_i < 0$ for exactly one
inequality $i \in \Active$, all other conditions holding exactly. Show that
dropping $i$ from the working set and re-solving~\eqref{eq:sub} strictly decreases
$f$. (This is the primal-method repair; the dual method of
Chapter~\ref{ch:walking} arranges never to be in this state.)
\end{exercise}

\begin{exercise}
The scaling of the tolerance $\varepsilon$ in~\eqref{eq:certificate} was left
vague. Propose a scaling that makes the test invariant under
$(G,a) \mapsto (\gamma G, \gamma a)$ for $\gamma > 0$, and another that makes it
invariant under $(c_i, b_i) \mapsto (\gamma c_i, \gamma b_i)$ per constraint.
Can both be achieved at once?
\end{exercise}

\begin{exercise}[Implementation]
Take the brute-force solver of Exercise~1.6 and replace its ``discard the
infeasible ones and compare objectives'' step by a certificate check: for each
candidate working set, form $(x,u)$ from~\eqref{eq:subsol} and test
\eqref{eq:certificate}. Confirm that exactly one working set passes on random
non-degenerate instances, and find an instance where more than one does. What is
special about it?
\end{exercise}

\end{exercises}
